% ============================================================
% Capitolo 6 — [MODULARE] Studio di caso: implementazione e risultati
%
% ------------------------------------------------------------
% ALTERNATIVE DI PROGETTO
%
% Opzione A — Analisi topologica descrittiva (rischio basso,
% consigliata).
%   Si implementa solo la costruzione del grafo dinamico e le sue
%   metriche (6.1, 6.2, 6.6), SENZA alcun modello predittivo. Si mostra
%   come densità, correlazione media e λ₂ evolvono attorno a eventi di
%   crisi noti.
%   Integrazione: il capitolo si riduce a ~7 pagine e cambia titolo in
%   "Analisi empirica della struttura di correlazione". Il Cap. 3 (GNN)
%   resta interamente teorico e diventa "background per sviluppi
%   futuri", da segnalare esplicitamente in 1.4. Vantaggio: non può
%   fallire — il grafo esiste e le sue metriche cambiano, punto. È un
%   risultato garantito.
%
% Opzione B — Confronto predittivo GCN vs. VAR vs. naive (rischio
% medio).
%   Il capitolo completo come descritto sotto (6.1-6.6).
%   Integrazione: richiede che il Cap. 4.5 (VAR) sia sviluppato con
%   cura e che il Cap. 5.4 (metriche) sia rigoroso, perché è lì che si
%   gioca la credibilità. Avvertenza forte: la probabilità che la GCN
%   batta un random walk sui rendimenti giornalieri cripto è bassa.
%   Bisogna quindi impostare la tesi fin dall'introduzione (1.4) come
%   "verifica di un'ipotesi", non come "dimostrazione di un vantaggio".
%   Un risultato negativo, ben documentato, è una tesi valida; un
%   risultato negativo presentato come fallimento è una tesi debole.
%
% ------------------------------------------------------------
% ============================================================
\chapter{Studio di caso: implementazione e risultati}
\label{ch:case-study}

\section{Dati e preprocessing}
\label{sec:data-preprocessing}
% (~2 pp.)
% Fonte (es. Binance/CoinGecko API), universo di asset (10-20 tra le
% principali cripto per capitalizzazione), granularità (giornaliera),
% periodo (deve includere almeno un evento di crisi), pulizia, calcolo
% dei log-rendimenti.

\section{Costruzione del grafo dinamico}
\label{sec:dynamic-graph-construction}
% (~2 pp.)
% Finestra rolling di 30/60 giorni, correlazione di Pearson,
% thresholding, giustificazione dei parametri.

\section{Architettura implementata}
\label{sec:implemented-architecture}
% (~2 pp.)
% GCN a 2 layer (PyTorch Geometric), task di previsione (regressione
% del rendimento a t+1 o classificazione binaria del segno), split
% walk-forward.

\section{Baseline di confronto}
\label{sec:comparison-baselines}
% (~1,5 pp.)
% VAR, e almeno un naive (random walk) — perché il naive è il vero
% avversario in finanza.

\section{Risultati predittivi}
\label{sec:predictive-results}
% (~2 pp.)
% Tabelle comparative, con onestà sui risultati: è del tutto legittimo
% (e apprezzato) concludere che il vantaggio è marginale o assente.

\section{Analisi dell'evoluzione topologica}
\label{sec:topological-evolution-analysis}
% (~2,5 pp.)
% La parte più promettente: serie temporale di densità del grafo,
% correlazione media, λ₂, e loro comportamento in prossimità di eventi
% noti. Questa analisi funziona anche se la previsione fallisce.
